\documentclass[12pt,a4paper,UTF8]{ctexart}
\usepackage[a4paper,top=22mm,bottom=22mm,left=22mm,right=22mm,headheight=15pt]{geometry}
\newcommand{\Needspace}[1]{\par\begingroup\skip0=\lastskip\vskip-\skip0\dimen0=#1\relax\advance\dimen0 by\skip0\vskip\dimen0\penalty-100\vskip-\dimen0\vskip\skip0\penalty9999\endgroup}
\usepackage{amsmath,amssymb,mathtools,bm}
\usepackage{graphicx,xcolor,booktabs,longtable,array,enumitem,fancyhdr,caption,float,microtype,setspace}
\usepackage[hidelinks,unicode]{hyperref}
\usepackage[most]{tcolorbox}
\definecolor{ContestBlue}{HTML}{17365D}
\hypersetup{pdftitle={第079套},pdfauthor={物理竞赛教学资料}}
\setstretch{1.24}
\setlength{\parindent}{2em}
\setlength{\parskip}{0.35em}
\setlength{\emergencystretch}{2em}
\setlist{itemsep=0.35em,topsep=0.4em}
\captionsetup{font=small,labelfont=bf,justification=centering}
\setcounter{secnumdepth}{0}
\pagestyle{fancy}
\fancyhf{}
\fancyhead[L]{\small 第079套}
\fancyfoot[C]{\thepage}
\renewcommand{\headrulewidth}{0.4pt}
\ctexset{
  section={format=\Large\bfseries\color{ContestBlue},beforeskip=1.2em,afterskip=0.7em},
  subsection={format=\large\bfseries\color{ContestBlue},beforeskip=1.0em,afterskip=0.5em},
  subsubsection={format=\normalsize\bfseries,beforeskip=0.8em,afterskip=0.35em}
}
\usepackage{calc}
\newcounter{none}
\let\originalsubsection\subsection
\renewcommand{\subsection}{\Needspace{6\baselineskip}\originalsubsection}
\makeatletter
\newcommand{\tightlist}{\setlength{\itemsep}{0pt}\setlength{\parskip}{0pt}}
\newcommand{\maxwidth}{\ifdim\Gin@nat@width>\linewidth\linewidth\else\Gin@nat@width\fi}
\newcommand{\maxheight}{\ifdim\Gin@nat@height>0.72\textheight 0.72\textheight\else\Gin@nat@height\fi}
\newcommand{\pandocbounded}[1]{\begingroup\setkeys{Gin}{width=\maxwidth,height=\maxheight,keepaspectratio}#1\endgroup}
\makeatother
\begin{document}
\begin{titlepage}
  \centering
  \vspace*{58mm}
  {\Huge\bfseries\color{ContestBlue} 第079套\par}
  \vfill
  \vspace*{28mm}
\end{titlepage}
\section{第 1 题}\label{ux7b2c-1-ux9898}

两根完全相同的四分之一细圆弧由均匀金属制成，具有弹性。两根圆弧分别按图示方式将一端固定在墙上，另一端各悬挂一个质量相同的重物。形变很小，圆弧质量可以忽略。比较两种固定方式中重物下降的高度。

\begin{figure}[H]
\centering
\includegraphics[width=0.85\linewidth,height=\textheight,keepaspectratio,alt={图 1-1　两种四分之一圆弧的固定方式及相同的竖直载荷。}]{../figures/question-01.pdf}
\caption*{图 1-1　两种四分之一圆弧的固定方式及相同的竖直载荷。}
\end{figure}

\newpage

\section{第 2 题}\label{ux7b2c-2-ux9898}

一圆柱形电容器外半径为 \(R\)，内半径为 \(R-d\)，其中 \(d\ll R\)；长度为
\(l\)，总质量为
\(M\)。电容器由三根对称分布的绝缘细线悬挂，只能绕中心轴转动而不能水平移动。空间中有竖直向下、磁感应强度为
\(B\) 的匀强磁场。电容器已经充电，两极板间电势差为 \(V\)。

分别分析以下过程产生的转动：

\begin{enumerate}
\def\labelenumi{\arabic{enumi}.}
\tightlist
\item
  保持磁场不变，用位于端面的径向导线组将两极板连接，使电容器放电；
\item
  保持电容器带电，令磁场突然减小到零。
\end{enumerate}

\begin{figure}[H]
\centering
\includegraphics[width=0.846\linewidth,height=\textheight,keepaspectratio,alt={图 2-1　轴向磁场中的薄间隙圆柱电容器及其悬挂方式。}]{../figures/question-02.pdf}
\caption*{图 2-1　轴向磁场中的薄间隙圆柱电容器及其悬挂方式。}
\end{figure}

\newpage

\section{第 3 题}\label{ux7b2c-3-ux9898}

质量为 \(m\)、半径为 \(R\)
的乒乓球斜向下撞击水平桌面。碰撞前速度的水平分量为 \(v_0>0\)，角速度为
\(\omega_0\)；转轴水平且垂直于速度的水平分量。规定上旋
\(\omega>0\)、下旋 \(\omega<0\)，反弹后水平速度与 \(v_0\) 同向为正。

作如下假设：竖直方向碰撞完全弹性；接触期间的滑动距离可忽略；离开桌面瞬间接触点无滑动；忽略空气阻力。已知球绕相应直径的转动惯量为

\[I=\frac23mR^2.\]

分别讨论 \(\omega_0=0\)、\(\omega_0>0\)、\(\omega_0<0\) 时反弹瞬间的
\(v_f\)、\(\omega_f\) 及符号条件。

\begin{figure}[H]
\centering
\includegraphics[width=0.9\linewidth,height=\textheight,keepaspectratio,alt={图 3-1　乒乓球碰撞前后的水平速度和自转方向。}]{../figures/question-03.pdf}
\caption*{图 3-1　乒乓球碰撞前后的水平速度和自转方向。}
\end{figure}

\newpage

\section{第 4 题}\label{ux7b2c-4-ux9898}

一块均匀薄木板原本为直板，现弯成半圆形并夹在两个静止物块之间，与两物块均为点接触。木板中性层长度为
\(\pi R\)，厚度为 \(d\)，垂直纸面方向宽度为 \(a\)，杨氏模量为
\(Y\)；木板不接触地面。求每个物块受到木板的水平作用力 \(F\)。

\begin{figure}[H]
\centering
\includegraphics[width=0.794\linewidth,height=\textheight,keepaspectratio,alt={图 4-1　弯成半圆形的薄木板及两端点接触。}]{../figures/question-04.pdf}
\caption*{图 4-1　弯成半圆形的薄木板及两端点接触。}
\end{figure}

\newpage

\section{第 5 题}\label{ux7b2c-5-ux9898}

竖直平行板电容器的两极板固定，间距为 \(d\)、高度为 \(h\)、面积
\(A\gg d^2\)。忽略边缘效应和空气阻力，考虑重力。

一质量为 \(M\)、电荷量为 \(q\) 的小金属球用长为 \(L\)
的绝缘细线悬挂。细线悬点固定；电容器不带电时，小球位于两板中央，高度为
\(h/2\)。

\begin{enumerate}
\def\labelenumi{\arabic{enumi}.}
\tightlist
\item
  两板电势差恒为 \(V_0\) 时，求平衡角 \(\theta_0\)；将小球从略大于
  \(\theta_0\) 的角度由静止释放，证明其作小振动并求周期。
\item
  小球通过平衡位置时剪断细线，求小球不撞击极板时 \(V_0\) 的最大值。
\item
  另一过程中，\(t=0\)
  时小球位于两板中心且静止，同时剪断细线，并加交流电压
  \(V(t)=V_0\sin\omega t\)。分别在 \(g\gg h\omega^2\) 和
  \(g\ll h\omega^2\) 两个极限下求小球不碰极板时 \(\omega\) 的范围。
\end{enumerate}

\begin{figure}[H]
\centering
\includegraphics[width=0.528\linewidth,height=\textheight,keepaspectratio,alt={图 5-1　平行板间带电摆球的受力与几何位置。}]{../figures/question-05.pdf}
\caption*{图 5-1　平行板间带电摆球的受力与几何位置。}
\end{figure}

\newpage

\section{第 6 题}\label{ux7b2c-6-ux9898}

用性质均匀、单位长度电阻恒定的电阻丝构造一个无限递归三角网络。最外层是边长为
\(l\) 的正三角形；在每个角落继续生成边长为上一层 \(1/3\)
的相似三角形，三个角落完全相同。长度为 \(l\) 的电阻丝电阻为 \(R_0\)。

\begin{enumerate}
\def\labelenumi{\arabic{enumi}.}
\tightlist
\item
  求顶点 \(A\) 与左下顶点 \(B\) 间的等效电阻；
\item
  空间中有均匀变化的垂直磁场
  \(\partial B/\partial t=k_0\)。定义最外层三角形磁通变化率
\end{enumerate}

\[\mathcal E_0=k_0\frac{\sqrt3}{4}l^2.\]

取沿最外边 \(AB\) 的路径，求从 \(A\) 到 \(B\) 的电阻压降线积分
\(U_{AB}=\sum I_ir_i\)； 3. 求整个网络的焦耳热功率。

\begin{figure}[H]
\centering
\includegraphics[width=0.707\linewidth,height=\textheight,keepaspectratio,alt={图 6-1　三角形角落按比例 1/3 无限递归的电阻网络。}]{../figures/question-06.pdf}
\caption*{图 6-1　三角形角落按比例 \(1/3\) 无限递归的电阻网络。}
\end{figure}

\newpage

\section{第 7 题}\label{ux7b2c-7-ux9898}

\begin{enumerate}
\def\labelenumi{\arabic{enumi}.}
\tightlist
\item
  体积为 \(V\)
  的房间由绝热材料筑成，气密性良好，仅有一条微小缝隙与外界相通，外界压强恒定。空气视为摩尔质量
  \(M_a\)、绝热指数 \(\gamma=1.4\) 的理想气体；温度为 \(T_0\)
  时房内空气密度为 \(\rho\)。求把房间缓慢加热到 \(T_1\) 所需热量。
\item
  两个与外界绝热的物体初温分别为 \(T_a\)、\(T_b\)，且
  \(T_a>T_b\)；恒定热容分别为
  \(C_a\)、\(C_b\)。用可逆卡诺热机连接两物体并对外做功，直至达到共同温度。求末温和最大做功。
\end{enumerate}

\begin{figure}[H]
\centering
\includegraphics[width=0.892\linewidth,height=\textheight,keepaspectratio,alt={图 7-1　定压泄气房间与由可逆热机连接的有限热源、冷源。}]{../figures/question-07.pdf}
\caption*{图 7-1　定压泄气房间与由可逆热机连接的有限热源、冷源。}
\end{figure}

\newpage

\section{第 8 题}\label{ux7b2c-8-ux9898}

把每个原子视为可极化介质小球，其感生电偶极矩满足
\(\bm p=\alpha\bm E\)。一个双原子分子由两个质量均为 \(m\)
的小球和一根劲度系数为 \(k\)、原长为 \(R\)
的轻弹簧组成。忽略一个原子产生的电场对另一个原子极化率的反馈，且形变量远小于
\(R\)。

\begin{enumerate}
\def\labelenumi{\arabic{enumi}.}
\tightlist
\item
  当 \(k\to\infty\)，沿原子连线施加恒定外电场
  \(E\)，求两原子间电磁作用力大小。
\item
  \(k\) 很大但有限，沿原子连线施加
  \(E(t)=E_0\cos\omega t\)。求稳态时原子间距改变量 \(\Delta(t)\)
  的平均值、振幅和角频率。
\item
  若外电场垂直于原子连线，结果如何？
\end{enumerate}

\begin{figure}[H]
\centering
\includegraphics[width=0.9\linewidth,height=\textheight,keepaspectratio,alt={图 8-1　可极化双原子模型及连接两原子的弹簧。}]{../figures/question-08.pdf}
\caption*{图 8-1　可极化双原子模型及连接两原子的弹簧。}
\end{figure}
\end{document}
